\section*{Regional interannual consistency}

\subsection*{Sampling and temporal pairing}

The analysis compares consecutive-year changes in ChinaTherm30 and Landsat 8 annual LST composites during 2013--2024. The Landsat date window, temperature conversion and strict quality screening follow the main-text Landsat 8 comparison. Each year requires at least three valid observations; both annual ChinaTherm30 3$\times$3 means must be available, and neighbourhood water fractions must not exceed 10\%. The original 20,000 candidate points were supplemented by a pool of 59,626 candidates to improve representation where cloud screening had reduced temporal support. Within each 4° tile--macroregion--elevation stratum, eligible original points were retained first, followed by eligible reserve points in their preassigned order until the original stratum quota was reached.

For each of the 11 adjacent-year pairs, both datasets use exactly the same selected locations. A stratum is included only if it has at least three paired points and a Kish effective sample size of at least three in every year pair. The retained stratum set and its area weights remain fixed across all pairs; individual points may change between pairs. Original 2° sampling weights are normalized within each retained stratum to its original area. Supported area fractions are relative to the full macroregional sampling-design area, not the fraction of pixels directly observed.

\subsection*{Metrics and uncertainty}

For product $P$, the regional change ending in year $t$ is
\begin{equation}
\Delta T^P_{r,t}=\frac{\sum_{s\in S_r}A_s\left[\sum_{i\in I_{s,t}}w_i(T^P_{i,t}-T^P_{i,t-1})/\sum_{i\in I_{s,t}}w_i\right]}{\sum_{s\in S_r}A_s},
\end{equation}
where $S_r$ is the fixed set of supported strata in region $r$, $A_s$ is stratum area, and $I_{s,t}$ contains the paired locations. Pearson correlation and RMSE are calculated from the 11 regional change values for the two datasets. The standard-deviation ratio is $\operatorname{SD}(\Delta T^{\mathrm{CT30}})/\operatorname{SD}(\Delta T^{\mathrm{L8}})$; it describes variability in adjacent-year changes, not the variability or trend of annual temperature levels. Direction agreement counts year pairs for which both changes have the same sign.

Confidence intervals use 1,000 spatial-cluster bootstrap replicates over 4° tiles. They are conditional on the observed years and supported domain and do not quantify uncertainty over unobserved regions or resampled years. SWM contains only nine spatial blocks. Observation-date summaries first calculate the median valid-observation DOY for each point and year, then take an area-weighted median across paired points; they are not the first and last acquisition dates of a season.

\subsection*{Results}

Regional change correlations range from 0.399 in SWM to 0.870 in HSE, and change RMSE ranges from 1.017 to 1.752~$^{\circ}$C (Figure~S8; Tables~S4--S5). HSE agrees in change direction for all 11 pairs and TP for ten pairs; ANW and SWM agree for six. Standard-deviation ratios of 0.349--0.826 indicate smaller adjacent-year changes in ChinaTherm30 than in Landsat 8 across all six regions. Fixed strata represent more than 99\% of the design area in TP, NEC, NCLP and ANW, compared with 83.3\% in HSE and 62.5\% in SWM. Year-pair sample sizes, temperature changes and observation support are reported in Table~S6.

\begingroup\small\setlength{\tabcolsep}{4pt}
\begin{longtable}{@{}lrrrrrr@{}}
\caption{Regional interannual-change metrics for 2013--2024 (11 year pairs). Supported area denotes the fixed-stratum fraction of the design area. SD ratio is ChinaTherm30/Landsat 8.}\label{tab:s4-interannual-metrics}\\
\toprule
Region & Area (\%) & Paired points & $r$ & RMSE ($^{\circ}$C) & SD ratio & Same sign \\
\midrule
\endfirsthead
\toprule
Region & Area (\%) & Paired points & $r$ & RMSE ($^{\circ}$C) & SD ratio & Same sign \\
\midrule
\endhead
\bottomrule
\endfoot
TP & 99.2 & 3,841--4,228 & 0.734 & 1.017 & 0.826 & 10/11 \\
NEC & 99.9 & 3,689--3,691 & 0.592 & 1.158 & 0.481 & 8/11 \\
NCLP & 99.3 & 1,871--1,883 & 0.773 & 1.363 & 0.349 & 8/11 \\
ANW & 99.6 & 4,381--4,417 & 0.465 & 1.752 & 0.423 & 6/11 \\
HSE & 83.3 & 1,665--2,334 & 0.870 & 1.269 & 0.432 & 11/11 \\
SWM & 62.5 & 407--674 & 0.399 & 1.329 & 0.520 & 6/11 \\
\end{longtable}

\begin{longtable}{@{}lrrrr@{}}
\caption{Spatial-cluster bootstrap 95\% intervals for the metrics in Table~S4 (1,000 replicates).}\label{tab:s5-interannual-ci}\\
\toprule
Region & Tile blocks & $r$ & RMSE ($^{\circ}$C) & SD ratio \\
\midrule
\endfirsthead
\toprule
Region & Tile blocks & $r$ & RMSE ($^{\circ}$C) & SD ratio \\
\midrule
\endhead
\bottomrule
\endfoot
TP & 28 & [0.394, 0.874] & [0.857, 1.344] & [0.706, 1.113] \\
NEC & 24 & [0.213, 0.764] & [0.790, 2.000] & [0.311, 0.673] \\
NCLP & 14 & [0.496, 0.858] & [1.018, 2.062] & [0.272, 0.494] \\
ANW & 27 & [0.200, 0.689] & [1.445, 2.159] & [0.360, 0.529] \\
HSE & 15 & [0.759, 0.905] & [0.836, 1.904] & [0.325, 0.613] \\
SWM & 9 & [0.158, 0.475] & [1.226, 1.839] & [0.318, 0.655] \\
\end{longtable}

\begin{longtable}{@{}llrrrrr@{}}
\caption{Paired-year regional changes and observation support. Temperature changes are in degrees Celsius. Count and DOY columns give area-weighted medians for the earlier/later year, respectively; DOY summarizes point-level median observation dates.}\label{tab:s6-interannual-pairs}\\
\toprule
Region & Year pair & Points & $\Delta T_{\rm CT30}$ & $\Delta T_{\rm L8}$ & Count & DOY \\
\midrule
\endfirsthead
\toprule
Region & Year pair & Points & $\Delta T_{\rm CT30}$ & $\Delta T_{\rm L8}$ & Count & DOY \\
\midrule
\endhead
\bottomrule
\endfoot
TP & 2013--2014 & 4,080 & -0.938 & -2.330 & 5/5 & 224/238 \\
TP & 2014--2015 & 4,113 & 1.173 & 1.405 & 5/6 & 238/242 \\
TP & 2015--2016 & 4,165 & -1.161 & -1.624 & 6/5 & 242/238 \\
TP & 2016--2017 & 4,025 & -1.057 & -0.224 & 5/5 & 239/248 \\
TP & 2017--2018 & 3,841 & -0.305 & -0.627 & 5/5 & 247/239 \\
TP & 2018--2019 & 3,843 & 0.823 & 0.685 & 5/5 & 237/237 \\
TP & 2019--2020 & 4,022 & 2.021 & 0.377 & 5/6 & 237/250 \\
TP & 2020--2021 & 4,224 & -1.232 & -0.091 & 6/5 & 251/242 \\
TP & 2021--2022 & 4,158 & 1.014 & 2.598 & 6/5 & 242/221 \\
TP & 2022--2023 & 4,121 & -2.110 & -2.422 & 5/5 & 222/242 \\
TP & 2023--2024 & 4,228 & -0.581 & 0.741 & 5/6 & 242/235 \\
NEC & 2013--2014 & 3,691 & 1.566 & 2.327 & 6/7 & 246/238 \\
NEC & 2014--2015 & 3,691 & -0.048 & 0.889 & 7/8 & 238/233 \\
NEC & 2015--2016 & 3,691 & 0.196 & -0.185 & 8/7 & 233/234 \\
NEC & 2016--2017 & 3,691 & -0.647 & -0.144 & 7/8 & 234/234 \\
NEC & 2017--2018 & 3,691 & -0.257 & -2.139 & 8/7 & 236/246 \\
NEC & 2018--2019 & 3,691 & 0.071 & 1.295 & 7/7 & 246/250 \\
NEC & 2019--2020 & 3,691 & -0.405 & -0.255 & 7/6 & 250/232 \\
NEC & 2020--2021 & 3,691 & -0.295 & -2.504 & 7/6 & 232/249 \\
NEC & 2021--2022 & 3,689 & 0.895 & 1.545 & 6/7 & 250/248 \\
NEC & 2022--2023 & 3,689 & -0.400 & 0.901 & 7/8 & 248/249 \\
NEC & 2023--2024 & 3,691 & -0.926 & -0.007 & 8/7 & 249/242 \\
NCLP & 2013--2014 & 1,878 & 0.024 & -0.158 & 6/5 & 235/232 \\
NCLP & 2014--2015 & 1,877 & 0.445 & -0.085 & 6/7 & 232/239 \\
NCLP & 2015--2016 & 1,877 & -0.104 & 0.829 & 7/5 & 238/239 \\
NCLP & 2016--2017 & 1,871 & 0.029 & 1.579 & 5/5 & 238/216 \\
NCLP & 2017--2018 & 1,874 & -0.604 & -3.310 & 5/6 & 217/248 \\
NCLP & 2018--2019 & 1,883 & 1.203 & 2.760 & 6/6 & 249/229 \\
NCLP & 2019--2020 & 1,878 & -1.025 & -2.427 & 6/5 & 229/241 \\
NCLP & 2020--2021 & 1,877 & 0.135 & 2.351 & 5/6 & 241/227 \\
NCLP & 2021--2022 & 1,882 & 0.204 & 0.031 & 5/7 & 226/223 \\
NCLP & 2022--2023 & 1,883 & 0.128 & 0.577 & 6/6 & 224/235 \\
NCLP & 2023--2024 & 1,882 & -1.050 & -1.118 & 6/6 & 234/232 \\
ANW & 2013--2014 & 4,417 & -0.018 & -0.689 & 8/9 & 227/228 \\
ANW & 2014--2015 & 4,416 & -0.930 & -1.310 & 9/9 & 228/230 \\
ANW & 2015--2016 & 4,412 & 0.662 & 2.018 & 9/8 & 230/226 \\
ANW & 2016--2017 & 4,408 & 0.293 & 0.245 & 9/9 & 226/228 \\
ANW & 2017--2018 & 4,397 & 0.069 & -2.436 & 9/9 & 228/236 \\
ANW & 2018--2019 & 4,392 & -0.063 & 2.781 & 9/9 & 236/224 \\
ANW & 2019--2020 & 4,398 & 0.324 & -2.550 & 9/9 & 224/233 \\
ANW & 2020--2021 & 4,417 & 0.445 & 2.644 & 9/9 & 233/228 \\
ANW & 2021--2022 & 4,410 & 0.147 & -0.431 & 9/9 & 228/231 \\
ANW & 2022--2023 & 4,396 & -0.281 & 1.598 & 9/9 & 231/227 \\
ANW & 2023--2024 & 4,381 & -2.488 & -2.585 & 9/8 & 227/224 \\
HSE & 2013--2014 & 2,334 & -0.926 & -2.395 & 5/4 & 239/247 \\
HSE & 2014--2015 & 2,105 & -0.238 & -1.028 & 4/4 & 248/257 \\
HSE & 2015--2016 & 2,015 & 1.174 & 4.504 & 4/4 & 259/227 \\
HSE & 2016--2017 & 1,785 & -0.051 & -0.537 & 4/4 & 223/230 \\
HSE & 2017--2018 & 1,808 & -0.244 & -1.430 & 4/4 & 230/233 \\
HSE & 2018--2019 & 2,094 & 0.493 & 0.890 & 4/4 & 233/246 \\
HSE & 2019--2020 & 1,840 & -0.354 & -0.906 & 5/4 & 247/239 \\
HSE & 2020--2021 & 1,881 & 0.493 & 0.491 & 4/4 & 241/251 \\
HSE & 2021--2022 & 2,231 & 0.810 & 1.247 & 4/5 & 251/245 \\
HSE & 2022--2023 & 1,868 & -1.997 & -2.483 & 5/4 & 245/239 \\
HSE & 2023--2024 & 1,665 & 0.183 & 1.327 & 4/4 & 239/238 \\
SWM & 2013--2014 & 648 & -1.016 & 0.542 & 5/4 & 238/221 \\
SWM & 2014--2015 & 547 & 0.265 & -2.293 & 4/4 & 218/235 \\
SWM & 2015--2016 & 534 & 1.019 & 3.402 & 4/4 & 239/217 \\
SWM & 2016--2017 & 407 & -0.242 & 0.262 & 4/4 & 215/223 \\
SWM & 2017--2018 & 443 & -0.201 & -1.817 & 4/4 & 219/227 \\
SWM & 2018--2019 & 624 & 0.198 & 0.632 & 4/4 & 227/225 \\
SWM & 2019--2020 & 474 & -0.037 & -0.683 & 4/4 & 227/214 \\
SWM & 2020--2021 & 448 & -0.196 & 0.148 & 4/4 & 214/239 \\
SWM & 2021--2022 & 582 & 1.364 & 0.932 & 4/4 & 227/228 \\
SWM & 2022--2023 & 674 & -1.384 & -0.482 & 4/4 & 226/225 \\
SWM & 2023--2024 & 670 & -0.057 & 0.368 & 4/4 & 225/232 \\
\end{longtable}

\endgroup

